% SI Text 1: related literature; the long version's Section 2 with cross-references
% redirected to the PNAS main text and SI.
\section{Related literature}\label{sec:literature}

Our contribution sits at the intersection of four literatures: the empirical debate over concentrating or dispersing funds, the evidence on how well peer review predicts research outcomes, the case for funding lotteries, and formal models of science funding.

Whether funders should concentrate resources on a few researchers or spread them across many is a long-running dispute, and concentration is the current practice: the top 1\% of NIH principal investigators receive close to 10\% of research project grant funding, a share that has risen since 1985 \citep{KatzMatter2020, Lauer2021}. Empirical studies of funded portfolios generally find diminishing marginal returns to funding per researcher, and reviews of this literature conclude that the evidence favors more spreading than current practice \citep{Mongeon2016, Aagaard2020}. Concentration is also self-reinforcing: early grant winners accumulate advantages beyond the grant itself \citep{BolVaanRijt2018}. This dispute typically concerns the returns to funding: whether research output per researcher increases or decreases as grants concentrate. In our model, the answer also depends on the size of the funder's budget relative to researchers' total resources: with a budget large enough to relieve most bottlenecks, the optimal allocation itself reaches nearly every researcher, so little is lost by spreading funds widely; with a tight budget, effectiveness requires concentrating funds on the researchers with the largest capability-resource gaps (main text). The value of any targeting, in turn, depends on what the funder can know about researchers (main text).

How well peer review predicts research outcomes is contested. Among funded NIH grants, percentile scores predict subsequent productivity barely better than chance \citep{Fang2016}, and reviewers assessing the same application agree only weakly \citep{Pier2018}. Studies with other designs find real but modest predictive power: better scores predict more publications, citations, and patents, including under controls for investigator history and institution \citep{LiAgha2015}, and a natural experiment in funding finds that higher-scored proposals produce more than lower-scored ones \citep{ParkLeeKim2015}. Both sides of this dispute measure review's accuracy: how well scores predict outcomes among funded applicants. Our model instead measures review's value: the expected output that review adds by reducing the funder's uncertainty about the capability-resource gap. That value depends on the field; the same noisy signal can be worth a great deal in one field and nothing in another (main text), and no single number settles the question.

Doubts about review have made lotteries a live proposal. Advocates argue that randomizing among applications that pass a review threshold saves review costs, reduces bias, and honestly reflects review's limited ability to discriminate among proposals that pass a review threshold \citep{FangCasadevall2016, Avin2019, Roumbanis2019}, and several funders now run partial lotteries \citep{Liu2020, Heyard2022, Luebber2025}. The existing formal argument for lotteries measures what a lottery saves: \citet{GrossBergstrom2019} model grant competition as a contest in which scientists invest costly effort in proposals, show that when only a small fraction of proposals is funded this effort can rival the scientific value of the research funded, and show that a partial lottery among proposals above a threshold removes the link between the wasted effort and the fraction funded. What has not been measured is what a lottery gives up: the value of targeting. Our model measures that value. It is negligible where budgets are ample, review is uninformative, or capability is evenly spread, and substantial where capability is heavy-tailed and budgets are tight (main text). Together the two measures let a funder weigh the choice field by field.

Formal models of science funding have concentrated on the incentive costs of competition and on the value of exploration. Beyond the contest model above, \citet{GrossBergstrom2025} find, in a game-theoretic model of contests for scarce rewards, that intensifying competition pushes scientists toward higher-risk, higher-return projects. \citet{AvinBJPS2019} simulates funding strategies on an epistemic landscape and finds that allocation with an explicitly random component significantly outperforms a peer-review-like strategy on large landscapes, because review concentrates effort on regions already known to be productive. \citet{Sorzano2026} model allocation as a decision under heavy-tailed uncertainty about project value, and combine a citation-based signal with a lottery weighted by that signal. In each of these models, funding raises a single quantity, a project's chance of success or a researcher's productivity, and none treats the timing of spending as the funder's choice. In our model, a researcher's research output increases in both capability and resources but is limited by the scarcer of the two, and the funder observes neither directly. This formulation lets us address questions raised by all four literatures: how widely funds should be spread, what peer review is worth, and when lotteries lose little.
